\documentclass[11pt]{article}
\newcommand{\paso}{6}
\usepackage[T1]{fontenc}
\usepackage[utf8]{inputenc}
\usepackage[spanish,es-nodecimaldot,es-noquoting,es-lcroman]{babel}
\spanishplainpercent
\usepackage{lmodern}
\usepackage[a4paper,margin=21mm,headheight=15pt]{geometry}
\usepackage{amsmath,amssymb,bm}
\usepackage{booktabs,array}
\usepackage{tikz}
\usetikzlibrary{arrows.meta,positioning,calc}
\usepackage[most]{tcolorbox}
\usepackage{microtype}
\usepackage{fancyhdr}
\usepackage{hyperref}
\hypersetup{colorlinks=true,linkcolor=blue!45!black,urlcolor=blue!45!black}
\setlength{\parindent}{0pt}
\setlength{\parskip}{5pt}
\setlength{\emergencystretch}{2em}
\pagestyle{fancy}
\fancyhf{}
\fancyhead[L]{\small BEM del cilindro: paso \paso}
\fancyhead[R]{\small Guía matemática independiente}
\fancyfoot[C]{\thepage}
\newcommand{\vct}[1]{\bm{#1}}
\newcommand{\dd}{\,\mathrm{d}}
\newcommand{\ii}{\mathrm{i}}
\newcommand{\norm}[1]{\left\lVert#1\right\rVert}
\newcommand{\VP}{\operatorname{VP}}
\newcommand{\tr}{\operatorname{tr}}
\newtcolorbox{idea}[1]{breakable,colback=blue!3,colframe=blue!50!black,
  title={#1},fonttitle=\bfseries,boxrule=0.7pt}
\newtcolorbox{alerta}[1]{breakable,colback=orange!4,colframe=orange!65!black,
  title={#1},fonttitle=\bfseries,boxrule=0.7pt}
\newtcolorbox{resultado}[1]{breakable,colback=teal!3,colframe=teal!55!black,
  title={#1},fonttitle=\bfseries,boxrule=0.7pt}
\newcommand{\inicio}[2]{%
 {\small\color{blue!55!black}\bfseries SERIE: LOS SIETE PASOS MATEMÁTICOS DEL BEM}\par
 {\LARGE\bfseries Paso \paso. #1\par}\vspace{3mm}
 \begin{idea}{¿Qué busca este paso?}#2\end{idea}}
\newcommand{\mapa}[3]{%
 \begin{center}
 \begin{tikzpicture}[node distance=4mm,
   caja/.style={draw=blue!50!black,fill=blue!3,rounded corners=2pt,
     text width=13cm,align=center,inner sep=6pt},
   >={Stealth[length=2mm]}]
 \node[caja] (a){#1};
 \node[caja,below=of a] (b){#2};
 \node[caja,below=of b] (c){#3};
 \draw[->,thick,blue!60!black] (a)--(b);
 \draw[->,thick,blue!60!black] (b)--(c);
 \end{tikzpicture}
 \end{center}}
\newcommand{\base}{%
 \smallskip{\footnotesize\color{black!65}
 Modelo de referencia: cilindro sólido, material homogéneo e isótropo,
 deformación plana, pequeñas deformaciones, sin fuerzas de volumen;
 amplitudes complejas con \(e^{+\ii\omega t}\), \(\omega>0\).}\par}
\newcommand{\fuentes}{%
 \par\medskip\begingroup\footnotesize
 \textbf{Referencia y continuación.}
 Este documento desarrolla una etapa del esquema
 \texttt{esquema\_ecuaciones\_bem.tex}. Las ecuaciones y los ejemplos se
 contrastaron con \texttt{bem\_cilindro.py}, ambos en la carpeta superior.
 La guía general y sus deducciones ampliadas se encuentran en
 \texttt{../guia\_bem/guia\_bem.pdf}. No se modifica la implementación.
 \par Para recompilar este documento, mantener \texttt{estilo.tex} en
 la misma carpeta y ejecutar \texttt{pdflatex} sobre su archivo
 \texttt{.tex}, desde esa carpeta.\par\endgroup}

\begin{document}
\inicio{Ensamblar y regularizar}{Transformar la identidad integral continua
en un sistema finito de ecuaciones, sin perder el término libre ni tratar
incorrectamente la singularidad. El resultado son dos matrices globales
\(H(\omega)\) y \(G(\omega)\); la carga se aplicará en el paso 7.}
\base
\mapa{Identidad: \(c\,\vct u+\VP\int_\Gamma T\vct u\,\dd\Gamma
 =\int_\Gamma U\vct t\,\dd\Gamma\)}
{Interpolar campos, colocar la fuente en cada nodo e integrar cada elemento}
{Sistema global: \(H(\omega)\vct u=G(\omega)\vct t\)}

\section*{1. Convención para no mezclar tres tipos de índices}
\begin{center}
\begin{tabular}{cl}
\toprule
Índice & Función\\
\midrule
\(c\) & Nodo de colocación o fila global por bloques.\\
\(g\) & Nodo global al que pertenece una incógnita o dato.\\
\(a\) & Nodo local \(1,2,3\) dentro de un elemento \(e\).\\
\(i,j\) & Componentes cartesianas \(x,y\) del núcleo y de los campos.\\
\(q\) & Punto de cuadratura.\\
\bottomrule
\end{tabular}
\end{center}
El símbolo \(c(\vct\xi)\) del término libre es un tensor, no el índice
\(c\). Cuando haga falta distinguirlos se escribirá \(C_c\).

\section*{2. Colocación nodal}
Se fija la fuente en \(\vct\xi_c=\vct x_c\). Para \(i=x,y\),
\[
 (C_c)_{ij}u_{c,j}+
 \sum_e\int_{\Gamma_e}T_{ij}u_j\,\dd\Gamma
 =\sum_e\int_{\Gamma_e}U_{ij}t_j\,\dd\Gamma,
\]
entendiendo el valor principal donde corresponde.
Hay dos ecuaciones por cada nodo fuente, por lo que \(N\) nodos producen
\(2N\) ecuaciones escalares.

\begin{idea}{No se elige una ecuación por elemento}
Los elementos sirven para representar e integrar el contorno. Las ecuaciones
se imponen en los nodos de colocación. En este esquema hay \(2n_e\) nodos,
no \(n_e\) nodos de colocación.
\end{idea}

\newpage
\section*{3. Sacar los valores nodales fuera de las integrales}
Dentro del elemento \(e\):
\[
 u_j(\eta)\simeq\sum_{a=1}^3P_a(\eta)u_{g(e,a),j},\qquad
 t_j(\eta)\simeq\sum_{a=1}^3P_a(\eta)t_{g(e,a),j}.
\]
Al sustituir:
\[
\begin{aligned}
 (C_c)_{ij}u_{c,j}
 &+\sum_{e,a}\left[\VP\int_{-1}^{1}
 T_{ij}P_aJ\,\dd\eta\right]u_{g(e,a),j}\\
 &=\sum_{e,a}\left[\int_{-1}^{1}
 U_{ij}P_aJ\,\dd\eta\right]t_{g(e,a),j}.
\end{aligned}
\]
Cada valor nodal es constante respecto a \(\eta\), por eso sale de la
integral. Lo que queda dentro es calculable.

\section*{4. Cada influencia es un bloque, no un número}
\[
 G^e_{ca}=\int_{-1}^{1}UP_aJ\,\dd\eta,\qquad
 \widehat H^e_{ca}=\VP\int_{-1}^{1}TP_aJ\,\dd\eta.
\]
Ambos son matrices \(2\times2\). Por ejemplo:
\[
 G^e_{ca}
 \begin{pmatrix}t_{g,x}\\t_{g,y}\end{pmatrix}
 =
 \begin{pmatrix}
 G^{e,xx}_{ca}t_{g,x}+G^{e,xy}_{ca}t_{g,y}\\
 G^{e,yx}_{ca}t_{g,x}+G^{e,yy}_{ca}t_{g,y}
 \end{pmatrix}.
\]
Las filas son las dos ecuaciones de la fuente; las columnas, las dos
componentes del nodo de destino.

\section*{5. Sumar contribuciones compartidas}
Si un nodo \(g\) pertenece a dos elementos, ambos aportan a la misma columna
global:
\[
 G_{cg}=\sum_{\{e,a:\,g(e,a)=g\}}G^e_{ca}.
\]
La misma regla se aplica a \(H\), incluyendo el término libre en su diagonal.
No se reemplaza una contribución anterior por la última.

\begin{resultado}{Dónde se coloca el bloque}
Con numeración global desde cero y componentes intercaladas, el bloque
\((c,g)\) ocupa las filas \(2c,2c+1\) y columnas \(2g,2g+1\).
Con ocho elementos, cada matriz es \(32\times32\).
\end{resultado}

\newpage
\section*{6. Separar exactamente la dificultad singular}
Se escribe
\[
 T=T^s+(T-T^s),\qquad
 H=H^s+\Delta H(\omega).
\]
La matriz estática contiene
\[
 H^s_{cg}=C_c\delta_{cg}
 +\sum_{\{e,a:\,g(e,a)=g\}}\VP\int_{\Gamma_e}T^sP_a\,\dd\Gamma.
\]
Para \(g\ne c\) no hay término libre. Si el elemento contiene la fuente,
la función de forma del nodo distinto de ella se anula allí.
Su cero compensa el comportamiento \(1/r\), permitiendo integrar esos
aportes sin calcular directamente la diagonal singular.

\section*{7. Deducir la diagonal con una traslación rígida estática}
En el problema auxiliar estático, sea
\(\vct u_g=\vct a\) para todos los nodos:
\[
 \nabla\vct u=\vct0
 \ \Longrightarrow\ \boldsymbol\varepsilon=\vct0
 \ \Longrightarrow\ \boldsymbol\sigma=\vct0
 \ \Longrightarrow\ \vct t=\vct0.
\]
Como \(\sum_aP_a=1\), la interpolación representa exactamente ese campo.
La ecuación estática discreta da
\[
 \sum_gH^s_{cg}\vct a=\vct0.
\]
Debe valer para cualquier vector \(\vct a\), luego
\[
 \boxed{\sum_gH^s_{cg}=0_{2\times2},\qquad
 H^s_{cc}=-\sum_{g\ne c}H^s_{cg}.}
\]
Se calcula primero la parte fuera de la diagonal y después se reconstruye
el bloque diagonal con esta igualdad.

\begin{alerta}{Qué contiene realmente esa diagonal}
\(H^s_{cc}\) incluye el término libre \(C_c\) más el valor principal estático
ponderado. No es el valor de \(T^s\) en \(r=0\), no se fija en cero y
no se debe añadir otro \(\tfrac12 I\).
La suma es de bloques \(2\times2\), no de todas las entradas escalares.
\end{alerta}

\newpage
\section*{8. Recuperar la parte dinámica}
La corrección se integra para todos los nodos, también el de colocación:
\[
 \Delta H^e_{ca}=
 \int_{-1}^{1}(T-T^s)P_aJ\,\dd\eta
 \simeq\sum_qW_{qa}(T_q-T_q^s).
\]
De manera independiente,
\[
 G^e_{ca}=\int_{-1}^{1}UP_aJ\,\dd\eta
 \simeq\sum_qW_{qa}U_q.
\]
La singularidad logarítmica de \(U\) es integrable con la cuadratura adaptada;
no se usa una diagonal de movimiento rígido para \(G\).

\mapa{Una vez: geometría, pesos, \(T^s\), bloques no diagonales de \(H^s\)}
{Cerrar \(H^s_{cc}=-\sum_{g\ne c}H^s_{cg}\)}
{Cada \(\omega\): evaluar \(U,T\), integrar \(G,\Delta H\), formar \(H=H^s+\Delta H\)}

\section*{9. Por qué la condición rígida no se impone a \(H(\omega)\)}
Una traslación constante en el espacio, pero armónica en el tiempo, tiene
aceleración \(-\omega^2\vct a\). Para \(\omega\ne0\) no satisface la ecuación
sin fuerzas de volumen:
\[
 \nabla\!\cdot\boldsymbol\sigma+\rho\omega^2\vct a
 =\rho\omega^2\vct a\ne\vct0.
\]
Por eso \(\sum_gH_{cg}(\omega)\) no tiene por qué ser cero.
Reemplazar la diagonal dinámica por menos la suma de sus otros bloques
suprimiría parte de la física de inercia.

\begin{alerta}{Regularización algebraica y evaluación numérica}
La resta \(T-T^s\) es regular en el sentido de eliminar la parte \(1/r\),
pero evaluarla como diferencia de números grandes puede perder precisión.
El tratamiento singular correcto y la estabilidad de Bessel son dos
cuestiones distintas.
\end{alerta}

\newpage
\section*{10. La estructura final del sistema}
\[
 \begin{pmatrix}
 H_{00}&H_{01}&\cdots\\
 H_{10}&H_{11}&\cdots\\
 \vdots&\vdots&\ddots
 \end{pmatrix}
 \begin{pmatrix}\vct u_0\\\vct u_1\\\vdots\end{pmatrix}
 =
 \begin{pmatrix}
 G_{00}&G_{01}&\cdots\\
 G_{10}&G_{11}&\cdots\\
 \vdots&\vdots&\ddots
 \end{pmatrix}
 \begin{pmatrix}\vct t_0\\\vct t_1\\\vdots\end{pmatrix}.
\]
Cada entrada dibujada es un bloque \(2\times2\).
Las matrices son densas: un elemento lejano sigue influyendo en la fuente.
No tienen la conectividad local dispersa típica de una rigidez FEM.

\section*{11. Controles concretos antes de resolver}
\begin{itemize}
 \item \textbf{Tamaño:} \(H,G\in\mathbb C^{2N\times2N}\).
 \item \textbf{Unidades:}
 \([U]=\mathrm{Pa}^{-1}\), \([T]=\mathrm m^{-1}\), de donde
 \([G]=\mathrm{m/Pa}\) y \(H\) es adimensional.
 Entonces \(H\vct u\) y \(G\vct t\) tienen unidades de metros.
 \item \textbf{Identidad rígida:} verificar
 \(\max_c\norm{\sum_gH^s_{cg}}\) cerca del error de redondeo.
 Con ocho elementos, la máxima entrada absoluta observada de esa suma
 es aproximadamente \(1.2\times10^{-16}\).
 \item \textbf{Acumulación:} un nodo compartido debe recibir todos los
 aportes de sus elementos.
 \item \textbf{Conservación del precálculo:} la corrección dinámica debe
 sumarse a una copia de \(H^s\), no modificar permanentemente la matriz
 estática al cambiar la frecuencia.
\end{itemize}
\begin{idea}{Lo que un control algebraico no demuestra}
La suma estática cero se impone durante la construcción: cumplirla no prueba
que los núcleos ni las integrales fuera de la diagonal sean correctos.
Se necesita también una referencia física, refinamiento y controles de
cuadratura.
\end{idea}
\section*{12. Entrega al siguiente paso}
Ya se conoce el operador discreto del material y la geometría para una
frecuencia dada. El paso 7 forma el vector de presión, resuelve el sistema
y extrae la amplitud y fase radiales.
\fuentes
\end{document}
